\documentclass{beamer}
\usetheme{Madrid}
\usecolortheme{default}
\usepackage{listings}
\usepackage{xcolor}
\usepackage{ctex}
\setmonofont{DejaVu Sans Mono}

% Code listing settings
\lstset{
    language=Python,
    basicstyle=\footnotesize\ttfamily,
    keywordstyle=\color{blue},
    commentstyle=\color{green!60!black},
    stringstyle=\color{red},
    numbers=left,
    numberstyle=\footnotesize\color{gray},
    stepnumber=1,
    numbersep=5pt,
    backgroundcolor=\color{white},
    showspaces=false,
    showstringspaces=false,
    showtabs=false,
    frame=single,
    tabsize=2,
    breaklines=true,
    breakatwhitespace=false
}

\title{Lab4 Sharing}
\author{李鹏达}
\date{\today}

\begin{document}

\frame{\titlepage}

\section{Validity Properties}

\begin{frame}[allowframebreaks,fragile]{Validity Properties: Invariant Checking}
    \begin{itemize}
        \item \textbf{Invariant:} Binary Search Tree (BST) ordering.
        $$ \text{keys}(\text{left}) < \text{root\_key} < \text{keys}(\text{right}) $$
        \item \textbf{Test Goal:} Ensure that operations preserve the BST invariant.
        \item \textbf{Functions to Test (TODO 1):}
        \begin{enumerate}
            \item \texttt{find}: After locating a key, the result remains valid in a BST.
            \item \texttt{delete}: After merging two binary search trees, the result remains a valid binary search tree.
        \end{enumerate}
    \end{itemize}

    \framebreak

    \begin{lstlisting}
# TODO: After locating a key, the result remains valid in a BST.
@given(keys_strategy, trees_strategy)
def test_find_valid(key: int, bst: BST[int,int]) -> None:
    bst.find(key)
    assert(is_valid(bst))
    \end{lstlisting}

    \begin{lstlisting}
# TODO: After merging two binary search trees, the result remains a valid binary search tree.
@given(trees_strategy, trees_strategy)
def test_union_valid(bst1: BST[int,int], bst2: BST[int,int]) -> None:
    assert(is_valid(BST.union(bst1, bst2)))
    \end{lstlisting}

\end{frame}

\section{Postconditions Testing}

\begin{frame}[allowframebreaks,fragile]{Postconditions Testing: Result Verification}
    \begin{itemize}
        \item \textbf{Postcondition Idea:} Verify the expected state/result \textbf{after} an operation. \emph{Avoid replicating code logic in the test.}
        \item \textbf{Functions to Test (TODO 2):}
        \begin{enumerate}
            \item \texttt{delete}: Deletion must not affect the search results for \textbf{other} keys.
            \item \texttt{union}: Lookup in the merged BST must return the correct value, with keys from $\text{BST}_1$ having \textbf{precedence} over $\text{BST}_2$.
        \end{enumerate}
    \end{itemize}

        \framebreak

    \begin{lstlisting}
# TODO: Deletion should not affect the search results for other keys.
@given(keys_strategy, trees_strategy, keys_strategy)
def test_delete_post(key: int, bst: BST[int,int], search_key: int) -> None:
    found = bst.delete(key).find(search_key)
    expected = bst.find(search_key) if search_key != key else None
    assert found == expected
    \end{lstlisting}

    \framebreak

    \begin{lstlisting}
# TODO: After merging, the lookup should return the correct value, with keys from bst1 taking precedence over keys from bst2.
@given(trees_strategy, trees_strategy, keys_strategy)
def test_union_post(bst1: BST[int,int], bst2: BST[int,int], search_key:int) -> None:
    union_bst: BST[int,int] = BST.union(bst1, bst2)
    found = union_bst.find(search_key)
    bst1_found = bst1.find(search_key)
    expected = bst1_found if bst1_found else bst2.find(search_key)
    assert found == expected
    \end{lstlisting}
\end{frame}

\section{Metamorphic Testing}

\begin{frame}[allowframebreaks,fragile]{Metamorphic Testing: Relation Between Calls}
    \begin{itemize}
        \item \textbf{Metamorphic Relation (MR) Idea:} Check the relationship between the results of two \textbf{related} function calls: $\text{result}_1 \sim \text{result}_2$.
        \item \textbf{Functions to Test (TODO 3):}
        \begin{enumerate}
            \item \texttt{delete}: Establish a relationship by adding an \texttt{insert} operation (e.g., $\text{delete}.\text{insert} \sim \text{insert}.\text{delete}$).
            \item \texttt{union}: Establish a relationship by adding an \texttt{insert} operation.
        \end{enumerate}
        \item \textbf{Equivalence:} Use $\texttt{equivalent}(\text{bst}_1, \text{bst}_2)$ to check if key-value sets are equal, \textbf{disregarding structure}.
    \end{itemize}

        \framebreak

    \begin{lstlisting}
# TODO: establish a relationship between `delete.insert` and `insert.delete` by adding an insert. That is, by inserting before and after the deletion, one can determine whether the tree structures are equivalent.
# NOTE: Take care to assess the scenario where key1 equals key2.
@given(keys_strategy, keys_strategy, st.integers(), trees_strategy)
def test_delete_metamorph_by_insert(key1: int, key2: int, value: int, bst: BST[int,int]) -> None:
    deleted = bst.delete(key1).insert(key2, value)
    expected = bst.insert(key2, value).delete(key1)
    assert equivalent(deleted, expected)
    \end{lstlisting}

    \begin{lstlisting}
# TODO: To verify the outcome of the union operation, one may establish a relationship by adding an insert.
# NOTE: Please note that we have not introduced bugs into the union, but the metamorphic properties of this union are prone to errors in implementation. Correct properties will not detect bugs.
@given (keys_strategy, st.integers(), trees_strategy, trees_strategy)
def test_union_metamorph_by_insert(key: int, value: int, bst1: BST[int,int], bst2: BST[int,int]) -> None:
    unioned = BST.union(bst1, bst2).insert(key, value)
    expected = BST.union(bst1.insert(key, value), bst2)
    assert equivalent(unioned, expected)
    \end{lstlisting}
\end{frame}

\section{Model-based Properties Testing}

\begin{frame}[allowframebreaks,fragile]{Model-based Properties: Abstraction}
    \begin{itemize}
        \item \textbf{Model Idea:} Relate the concrete data type (BST) to an abstract data model (List of $\langle K, V \rangle$ tuples) via an $\mathbf{Abstraction}$ \textbf{Function} ($\texttt{to\_list()}$).
        \item \textbf{Test Goal:} $\texttt{concrete\_op}(\text{BST}) \sim \texttt{abstract\_op}(\text{model})$.
        \item \textbf{Functions to Test (TODO 4):}
        \begin{enumerate}
            \item \texttt{delete}: Compare the set of tuples after deleting from the BST versus deleting from the abstract list.
            \item \texttt{union}: Compare the set of tuples after performing \texttt{union} on two BSTs versus two abstract lists.
        \end{enumerate}
    \end{itemize}

            \framebreak

    \begin{lstlisting}
# TODO:Perform a delete operation on the BST and abstract data structure, then determine whether the final sets are equivalent.
@given(keys_strategy, trees_strategy)
def test_delete_model(key: int, bst: BST[int,int]) -> None:
    abstract_data = bst.to_list()
    deleted_abstract_data = filter(lambda kv: kv[0] != key, abstract_data)
    deleted_bst = bst.delete(key)
    assert set(deleted_bst.to_list()) == set(deleted_abstract_data)
    \end{lstlisting}

    \framebreak

  \begin{enumerate}
    \item Get keys in BST1
    \item Get key-value pairs in BST2 whose keys are not in BST1
    \item Merge the key-value pairs
  \end{enumerate}

    \begin{lstlisting}
# TODO:Perform a union operation on two BSTs and their corresponding abstract data structures, then determine whether the resulting union is equivalent.
@given(trees_strategy, trees_strategy)
def test_union_model(bst1: BST[int,int], bst2: BST[int,int]) -> None:
    abstract_data1 = bst1.to_list()
    abstract_data2 = bst2.to_list()
    data1_key_set = set(map(lambda kv: kv[0], abstract_data1))
    unioned_abstract_data = list(filter(lambda kv: kv[0] not in data1_key_set, abstract_data2)) + abstract_data1
    unioned_bst = BST.union(bst1, bst2)
    assert set(unioned_bst.to_list()) == set(unioned_abstract_data)
    \end{lstlisting}
\end{frame}

\begin{frame}
\begin{center}
    \huge THANKS
\end{center}
\end{frame}

\end{document}